% ============================================================
% CHAPTER 6: FEEDBACK VERTEX SET (matches Vazirani Ch. 6)
% ============================================================
\chapter{Feedback Vertex Set}
\label{ch:fvs}

\section{Intuition: Peeling and Layering}

\begin{intuitionbox}
\textbf{Peeling and Layering Intuition:} Feedback Vertex Set (FVS) is the problem of breaking all cycles in a graph by deleting a minimum-weight set of vertices. Low-degree vertices cannot support cycles: degree 0 or 1 vertices are useless for cycles, and degree 2 vertices can be contracted. In a ``clean'' graph (minimum degree $\ge 3$), the average degree of any forest is less than 2, meaning any cycle-free subgraph must connect heavily to the deleted feedback vertices. By reducing weights proportionally to a vertex's cycle-breaking contribution (measured by its cyclomatic weight $\delta_G(v)$), we pay at most twice the optimal cost to hit every cycle.
\end{intuitionbox}

\section{Problem Definition}
Given an undirected graph $G = (V, E)$ with non-negative vertex weights $w: V \to \mathbb{R}^+$, the \emph{Minimum Feedback Vertex Set} (FVS) problem is to find a subset $F \subseteq V$ of minimum total weight $\sum_{v \in F} w(v)$ such that the induced subgraph $G[V \setminus F]$ contains no cycles (i.e., it is a forest).

The decision version of FVS is one of Karp's 21 NP-complete problems~\cite{karp72}. It has applications in database deadlock resolution, circuit design, and Bayesian network inference (loop cutsets).

\section{Algorithm and Python Implementation}
The 2-approximation algorithm uses the local-ratio technique. It first simplifies the graph by removing vertices of degree $\le 1$ and contracting vertices of degree 2. It then scales weights based on the cyclomatic weight $\delta_H(v) = \text{cyc}(H) - \text{cyc}(H \setminus \{v\})$, where $\text{cyc}(H) = |E(H)| - |V(H)| + \text{cc}(H)$ is the cyclomatic number of graph $H$.

\begin{lstlisting}[style=python, caption=Local Ratio Algorithm for Feedback Vertex Set]
def connected_components_count(graph):
    """Number of connected components of the multigraph."""
    visited, count = set(), 0
    for v in graph:
        if v in visited:
            continue
        count += 1
        q = [v]
        visited.add(v)
        while q:
            curr = q.pop(0)
            for nxt in graph[curr]:
                if nxt in graph and nxt not in visited:
                    visited.add(nxt)
                    q.append(nxt)
    return count


def cyclomatic_number(graph):
    """Cyclomatic number m - n + cc of the multigraph."""
    n = len(graph)
    if n == 0:
        return 0
    m = sum(len(graph[v]) for v in graph) // 2
    return m - n + connected_components_count(graph)


def is_feedback_vertex_set(graph, fvs):
    """Is graph - fvs acyclic (a forest)?"""
    remaining = set(graph) - fvs
    visited = set()
    for start in remaining:
        if start in visited:
            continue
        q = [(start, -1)]
        visited.add(start)
        while q:
            curr, parent = q.pop(0)
            neighbors = [x for x in graph[curr] if x in remaining]
            if len(neighbors) != len(set(neighbors)):
                return False          # 2-cycle
            for nxt in neighbors:
                if nxt == parent:
                    continue
                if nxt in visited:
                    return False      # cycle
                visited.add(nxt)
                q.append((nxt, curr))
    return True


def feedback_vertex_set_approx(graph, weights):
    """2-approx for undirected weighted Feedback Vertex Set."""
    g = {v: list(neighbors) for v, neighbors in graph.items()}
    
    # 0. If any vertex has weight <= 0, remove and recurse
    for v in list(g.keys()):
        if weights.get(v, 0.0) <= 1e-9:
            g_sub = {x: [y for y in neighbors if y != v] for x, neighbors in g.items() if x != v}
            sub_fvs = feedback_vertex_set_approx(g_sub, weights)
            sub_fvs.add(v)
            # Prune before returning
            final_fvs = set(sub_fvs)
            for x in list(final_fvs):
                if is_feedback_vertex_set(graph, final_fvs - {x}):
                    final_fvs.remove(x)
            return final_fvs
            
    # 1. Clean the graph (remove degree <= 1, contract degree 2)
    changed = True
    while changed:
        changed = False
        to_remove = [v for v in g if len(g[v]) <= 1]
        if to_remove:
            for v in to_remove:
                for nxt in g[v]:
                    if nxt in g:
                        while v in g[nxt]: g[nxt].remove(v)
                del g[v]
            changed = True
            continue
            
        deg2_vertices = [v for v in g if len(g[v]) == 2]
        for v in deg2_vertices:
            if v not in g: continue
            u, w = g[v]
            if u != w:
                del g[v]
                g[u] = [w if x == v else x for x in g[u]]
                g[w] = [u if x == v else x for x in g[w]]
                changed = True
                break

    if not g or cyclomatic_number(g) == 0:
        return set()

    # 2. Check for self-loops (cycle of length 1)
    for v in g:
        if v in g[v]:
            g_sub = {x: [y for y in neighbors if y != v] for x, neighbors in g.items() if x != v}
            sub_fvs = feedback_vertex_set_approx(g_sub, weights)
            sub_fvs.add(v)
            return sub_fvs

    # 3. Check for multi-edges (cycle of length 2)
    for u in g:
        for v in g[u]:
            if g[u].count(v) > 1:
                eps = min(weights[u], weights[v])
                weights_next = dict(weights)
                weights_next[u] -= eps
                weights_next[v] -= eps
                g_sub = {x: list(neighbors) for x, neighbors in g.items()}
                sub_fvs = feedback_vertex_set_approx(g_sub, weights_next)
                if u not in sub_fvs and v not in sub_fvs:
                    sub_fvs.add(u if weights_next[u] <= 1e-9 else v)
                return sub_fvs

    # 4. Normal local ratio step
    cyc_g = cyclomatic_number(g)
    deltas = {}
    for v in g:
        g_without_v = {x: [y for y in neighbors if y != v] for x, neighbors in g.items() if x != v}
        deltas[v] = cyc_g - cyclomatic_number(g_without_v)
        if deltas[v] <= 0: deltas[v] = 1

    eps = min(weights[v] / deltas[v] for v in g)
    weights_next = dict(weights)
    for v in g:
        weights_next[v] -= eps * deltas[v]
        
    sub_fvs = feedback_vertex_set_approx(g, weights_next)
    for v in [v for v in g if weights_next[v] <= 1e-9]:
        if v not in sub_fvs: sub_fvs.add(v)
            
    final_fvs = set(sub_fvs)
    for v in list(final_fvs):
        if is_feedback_vertex_set(graph, final_fvs - {v}):
            final_fvs.remove(v)
    return final_fvs
\end{lstlisting}

\section{Analysis: Local Ratio and Cyclomatic Weights}
The cyclomatic number $\text{cyc}(G) = |E| - |V| + \text{cc}(G)$ represents the maximum number of independent cycles in $G$.
The value $\delta_G(v) = \text{cyc}(G) - \text{cyc}(G \setminus \{v\})$ measures how many cycles are broken by deleting $v$. We have:
\[
\delta_G(v) = d_G(v) - 1 - (k(G \setminus \{v\}) - k(G))
\]
Since a minimal feedback vertex set $F$ leaves a forest $T = G[V \setminus F]$, the average degree in $T$ is less than 2. From this, we can show that the sum of cyclomatic weights of any minimal FVS is bounded:
\[
\sum_{v \in F} \delta_G(v) \le 2 \cdot \text{OPT}_{IP}
\]
By the Local Ratio Theorem, decomposing the weights $w(v)$ into $w'(v) + \epsilon \cdot \delta_G(v)$ guarantees that the final minimal FVS satisfies $w(F) \le 2 \cdot \text{OPT}$.

\section{Applications}
\begin{itemize}
    \item \textbf{Bayesian Inference (Loop Cutsets):} Pearl's belief propagation algorithm is exact on trees. For graphs with cycles, belief propagation can be run by conditioning on a loop cutset (which is an FVS). Finding a minimum-weight FVS minimizes the conditioning overhead.
    \item \textbf{Deadlock Resolution:} In database transactions, the wait-for graph tracks transaction dependencies. Cycles represent deadlocks. An FVS represents a set of transactions to abort to resolve all deadlocks with minimal cost.
    \item \textbf{Dependency Resolution in Compilers:} Breaking cyclic package dependencies in package managers or modular architectures.
\end{itemize}

\subsection*{Real Example: Deadlock Resolution in Distributed Databases}
\textbf{Scenario:} A database has $n=100$ concurrent transactions. Cyclic dependencies (deadlocks) occur when transaction $A$ waits for $B$, $B$ waits for $C$, and $C$ waits for $A$. Aborting a transaction has a weight (representing processing time lost). We want to break all deadlocks with minimal total weight aborted.
\textbf{Model:} Directed graph of dependencies. For undirected cycles (simplifying directed FVS to undirected FVS when cycles are symmetric), the FVS algorithm finds the minimum-cost set of transactions to abort.
\textbf{Why 2-approx matters:} Minimizing transaction abort cost is NP-hard. A 2-approximation guarantees that we abort transactions costing no more than $2\times$ the absolute minimum abort cost required to free the database from deadlocks.

\section{Tight Example: Bipartite Graph}
For a bipartite graph $K_{3,3}$ where all vertex weights are 1, the optimal FVS is of size 2. If the algorithm selects a suboptimal set or order, the local ratio step can yield a set of size 4 if vertices are not chosen carefully. The ratio approaches 2 for larger instances.

\section{Exercises}
\begin{exercise}
Prove that for any tree $T$, the Feedback Vertex Set is empty.
\end{exercise}
\begin{exercise}
Prove that for any graph $G$, the average degree of a forest $T \subset G$ is less than 2.
\end{exercise}

